\documentclass[11pt]{article}
\usepackage[margin=1in]{geometry}
\usepackage{amsmath,amssymb,amsthm,booktabs,hyperref}
\newtheorem{theorem}{Theorem}
\newtheorem{lemma}{Lemma}
\newtheorem{proposition}{Proposition}
\newtheorem{corollary}{Corollary}
\title{An obstruction for symmetrized matrix multiplication sectors}
\author{AI-generated research draft prepared with Codex}
\date{October 10, 2026}
\begin{document}
\maketitle
\paragraph{Scope and attribution.}
This is an independent follow-up to
\href{https://github.com/openai/math/tree/fd4aeeb2ee4fc729c18d98444fed42fd0529eeeb/preprints/Matrix-Multiplication-Nine-Fourths-October-2-2026}{OpenAI's original manuscript},
pinned at commit \texttt{fd4aeeb2ee4fc729c18d98444fed42fd0529eeeb}.
It is not an official OpenAI publication. Mathematical priority and
independent formal verification are not claimed. Exact checks supplement
the proofs below and do not verify the imported source lemmas.
\section{An obstruction for the matrix profile inequalities}
The matrix paper defines a symmetric positive profile $P(a,b)$ with
$P(1,b)=b$, separate discrete concavity, and shifted tripling
$P(a,3h+a-1)\ge3P(a,h)$. Its rank envelope is
$P(a,b)\le(a+b-1)^{1/t}$. We exhibit an abstract profile compatible with
these conditions at $t=3/4$.

\begin{theorem}\label{thm:barrier}
For real $a,b\ge1$, set $u=\min(a,b)$, $v=\max(a,b)$ and
\[
 P_*(a,b)=\left(\frac{3u-1}{2}\right)^{1/3}
                      \left(v+\frac{u-1}{2}\right).
\]
Then $P_*$ is positive, symmetric, nondecreasing and separately concave,
$P_*(1,b)=b$, and it satisfies shifted tripling on positive integers.
Moreover
\[
 P_*(a,b)\le(a+b-1)^{4/3},\qquad
 P_*(a,a)=\left(\frac{3a-1}{2}\right)^{4/3}=\Theta(a^{4/3}).
\]
Thus these abstract constraints alone cannot force a larger diagonal
growth exponent, or exclude $t=3/4$.
\end{theorem}
\begin{proof}
Fix $a$. On $b\ge a$ the function is linear with positive slope.
On $1\le b\le a$, put $x=3b-1$, $A=3a-1$; it is
$(x/2)^{1/3}(2A+x)/6$. Its derivative in $b$ is positive and its
second derivative is
\[
 \frac{2^{2/3}(x-A)}{3x^{5/3}}\le0.
\]
The one-sided first derivatives match at $b=a$, both equal to
$((3a-1)/2)^{1/3}$. Concavity follows globally, and implies the required
discrete inequalities. Symmetry supplies the other coordinate.

If $h\ge a$, substitution gives
$P_*(a,3h+a-1)=3P_*(a,h)$ exactly. If $h<a$, put
$A=3a-1$, $B=3h-1$, and $z=(A/B)^{1/3}>1$.
After multiplication by a positive factor, shifted tripling becomes
\[
 z^4-2z^3+2z-1=(z-1)^3(z+1)\ge0.
\]
The boundary value and diagonal formula are direct substitutions.
Finally both $(3u-1)/2$ and $v+(u-1)/2$ are at most $a+b-1$,
proving the rank-envelope inequality.
\end{proof}
This is not a construction of a tensor character. It is an obstruction
to a proof using only the listed inequalities. New valid tensor
constraints, or a different construction, could still improve $9/4$.
The symbolic identity and exact rational cube comparisons were checked
by \texttt{tests/test\_profile.py}; these supplement the displayed proof.

\section{Matrix multiplication: a barrier for arbitrary matched sectors}
\label{sec:sector-barrier}
The earlier profile obstruction extends beyond three equal branches.
Consider a degeneration of $C(a,B)$ into matched blocks sharing the
first leg, each individually equivalent to $C(a,h_i)$, allowing an
exchange of its second and third legs. No auxiliary tensor factors
are allowed in the output in this statement. Suppose there are $s$
blocks, of which $r$ have exchanged legs. Their second and third spaces
are respectively disjoint. Counting their dimensions forces
\[
 B\ge\sum_i h_i+r(a-1),\qquad
 B+a-1\ge\sum_i h_i+(s-r)(a-1),
\]
and hence
\begin{equation}\label{eq:packing-bound}
 B\ge\sum_i h_i+\lfloor s/2\rfloor(a-1).
\end{equation}
Indeed $\min_{0\le r\le s}\max(r,s-r-1)=\lfloor s/2\rfloor$.
This constraint holds regardless of which substitutions produce the
matched blocks.

Applying the source's shared-leg entropy inequality with the same
probability vector to all six permuted characters, then optimizing it,
gives the symmetrized consequence
\[
 P(a,B)\ge\sum_{i=1}^sP(a,h_i).
\]
Each block's possible leg exchange merely permutes the six factors.

\begin{proposition}
The abstract profile $P_*$ of Theorem~\ref{thm:barrier} satisfies every
such consequence whenever the dimension bound~\eqref{eq:packing-bound}
holds. Thus allowing unequal lengths and arbitrarily many standard
matched sectors does not exclude this profile at $t=3/4$.
\end{proposition}
\begin{proof}
Put $d=(a-1)/2$ and $L=((3a-1)/2)^{1/3}$. We claim
$P_*(a,h)\le L(h+d)$ for every $h\ge1$, with equality for $h\ge a$.
For $h<a$, set $A=3a-1$, $H=3h-1$. Cubing both positive sides reduces
the claim to
\[
 A(A+2H)^3-H(2A+H)^3=(A-H)^3(A+H)\ge0.
\]
For $s\ge2$, the dimension bound gives $B\ge a$ and
\[
 P_*(a,B)=L(B+d)
 \ge L\left(\sum_i h_i+(2\lfloor s/2\rfloor+1)d\right)
 \ge\sum_iL(h_i+d)\ge\sum_iP_*(a,h_i).
\]
For $s=1$ use monotonicity. This proves the assertion.
\end{proof}
This rules out improving the exponent using just these symmetrized
matched-sector inequalities together with the earlier scalar constraints.
It leaves open more general output blocks, auxiliary factors, additional
tensor identities, and stronger information retained before symmetrization.
No improved matrix-multiplication exponent is claimed.


\section{Exact saturation and parity slack}
The sector barrier also has an explicit equality classification.
Let $s\ge2$, $a\ge1$, $d=(a-1)/2$, $L=((3a-1)/2)^{1/3}$ and
$B_{\min}=\sum_i h_i+2\lfloor s/2\rfloor d$. For $B\ge B_{\min}$,
we have $B\ge a$, hence
\begin{align*}
 P_*(a,B)-\sum_i P_*(a,h_i)
 &=L(B-B_{\min})
   +L(2\lfloor s/2\rfloor-s+1)d\\
 &\quad+\sum_i\bigl[L(h_i+d)-P_*(a,h_i)\bigr].
\end{align*}
All terms on the right are nonnegative. The earlier dominance identity
shows that each summand in brackets vanishes precisely when $h_i\ge a$:
it is strictly positive when $h_i<a$. For $a>1$, the middle term is zero
for odd $s$ and equals $Ld>0$ for even $s$. Thus odd sector counts attain
equality exactly when $B=B_{\min}$ and every $h_i\ge a$. Even counts
have gap at least $Ld$, with that lower bound attained under the same
conditions. If $a=1$, then $d=0$ and $P_*(1,h)=h$, so equality simply
requires $B=\sum_i h_i$ for either parity.

These are identities of the abstract profile and packing constraints.
They do not assert that every dimension-feasible packing is realizable
by a tensor degeneration. They identify equality families that the
retained inequalities already permit, and quantify the slack in the
remaining cases.

\section{The least admissible profile}\label{sec:least}
Call a function $P:\mathbb N_{>0}^2\to\mathbb R$ \emph{admissible} if it is symmetric,
satisfies $P(1,b)=b$, is discretely concave in each argument, and
satisfies shifted tripling $P(a,3h+a-1)\ge3P(a,h)$ for all $a,h\ge1$.
These are exactly the hypotheses of the source's diagonal-growth lemma;
positivity and the rank envelope are not used in this section.
Pointwise infima of admissible functions are admissible: concavity and
the boundary conditions are preserved trivially, and
$P_n(a,3h+a-1)\ge3P_n(a,h)\ge3\inf_mP_m(a,h)$ passes to the infimum
on the left. Hence the admissible functions have a least element.
We identify it in closed form.

For $u\ge1$ put
\[
 H_u=\prod_{m=1}^{u-1}\left(1+\frac1{3m}\right)
    =\frac{\Gamma(u+\tfrac13)}{\Gamma(u)\,\Gamma(\tfrac43)},
 \qquad
 P_{\min}(a,b)=H_u\left(v+\frac{u-1}{2}\right),
\]
with $u=\min(a,b)$, $v=\max(a,b)$ as before. Thus $P_{\min}$ has the
shape of $P_*$ with the cube root $((3u-1)/2)^{1/3}$ replaced by the
rational number $H_u$; the sequence $H_u$ is the product appearing in
the source's normalized-growth bound.

\begin{lemma}\label{lem:ratio}
For integers $1\le h\le a$,
\[
 H_a\,(a+2h-1)\ \ge\ H_h\,(2a+h-1),
\]
with equality if and only if $a\le h+1$.
\end{lemma}
\begin{proof}
Fix $h$ and set $g(a)=H_a(a+2h-1)-H_h(2a+h-1)$. Then $g(h)=0$, and
$g(h+1)=H_h\frac{3h+1}{3h}\cdot3h-H_h(3h+1)=0$. For $a\ge h+1$, using
$H_{a+1}=H_a\frac{3a+1}{3a}$,
\[
 g(a+1)-g(a)=H_a\Bigl(\tfrac{3a+1}{3a}(a+2h)-(a+2h-1)\Bigr)-2H_h
            =H_a\,\frac{4a+2h}{3a}-2H_h.
\]
If $g(a)\ge0$, then $H_a\ge H_h\frac{2a+h-1}{a+2h-1}$, and so
\[
 g(a+1)-g(a)\ \ge\ H_h\,
 \frac{(2a+h-1)(4a+2h)-6a(a+2h-1)}{3a(a+2h-1)}
 =H_h\,\frac{2(a-h)(a-h+1)}{3a(a+2h-1)},
\]
which is strictly positive for $a\ge h+1$. By induction $g(a)\ge0$ for
all $a\ge h$, with $g(a)>0$ as soon as $a\ge h+2$.
\end{proof}

\begin{theorem}\label{thm:least}
$P_{\min}$ is admissible, and every admissible $P$ satisfies
$P\ge P_{\min}$ pointwise. Moreover
\[
 P_{\min}(a,a)=\frac{3a-1}{2}\,H_a,\qquad
 \frac{P_{\min}(a,a)}{a^{4/3}}\longrightarrow\frac{3}{2\,\Gamma(4/3)}
 \approx1.67977,
\]
and $P_{\min}\le P_*\le(a+b-1)^{4/3}$, where the ratio
$P_*(a,b)/P_{\min}(a,b)=((3u-1)/2)^{1/3}/H_u$ increases strictly with
$u$ from $1$ to the limit $(3/2)^{1/3}\Gamma(4/3)\approx1.02221$.
\end{theorem}

\begin{proof}
\emph{Admissibility.} Symmetry and $P_{\min}(1,b)=b$ are immediate
($H_1=1$). Fix $a$ and write $\Delta_b=P_{\min}(a,b+1)-P_{\min}(a,b)$.
For $b\ge a$ the function is linear in $b$, so $\Delta_b=H_a$. For
$1\le b\le a-1$, symmetry gives $P_{\min}(a,b)=H_b(a+\frac{b-1}2)$, whence
\[
 \Delta_b=H_{b+1}(a-b)-H_b(a-b-1).
\]
At $b=a-1$ this is $H_a$, matching the linear part. For $b\le a-2$,
substituting $H_{b+1}=H_b\frac{3b+1}{3b}$ and
$H_{b+2}=H_{b+1}\frac{3b+4}{3b+3}$ gives
\[
 \Delta_b-\Delta_{b+1}=(a-b-1)\bigl(2H_{b+1}-H_b-H_{b+2}\bigr)
 =\frac{2(a-b-1)\,H_b}{9b(b+1)}\ \ge0 .
\]
So the increments are nonincreasing in $b$: $P_{\min}$ is discretely
concave in each argument.

For tripling, note $3h+a-1\ge a$, so
$P_{\min}(a,3h+a-1)=\tfrac32H_a(a+2h-1)$. If $h\ge a$ then
$3P_{\min}(a,h)=\tfrac32H_a(a+2h-1)$ and tripling holds with equality.
If $h<a$ then $3P_{\min}(a,h)=\tfrac32H_h(2a+h-1)$, and tripling is
exactly Lemma~\ref{lem:ratio}, with equality precisely when $h=a-1$.

\emph{Minimality.} Let $P$ be admissible, fix $a$, and write
$\Delta_{a,h}=P(a,h+1)-P(a,h)$; by concavity these are nonincreasing in
$h$. We first show
\begin{equation}\label{eq:slope}
 \Delta_{a,h}\ \ge\ \frac{2P(a,a)}{3a-1}\qquad(h\ge a).
\end{equation}
Indeed, tripling and monotonicity of the increments give
\[
 3P(a,h)\le P(a,3h+a-1)=P(a,h)+\sum_{j=h}^{3h+a-2}\Delta_{a,j}
 \le P(a,h)+(2h+a-1)\Delta_{a,h},
\]
while $P(a,h)=P(a,a)+\sum_{j=a}^{h-1}\Delta_{a,j}\ge P(a,a)+(h-a)\Delta_{a,h}$.
Combining, $2P(a,a)+2(h-a)\Delta_{a,h}\le(2h+a-1)\Delta_{a,h}$, which
is~\eqref{eq:slope}.

Put $H^P_a=2P(a,a)/(3a-1)$, so $H^P_1=1$. For $a\ge2$, symmetry,
\eqref{eq:slope} applied at $(a-1,a-1)$, and concavity give
\[
 P(a,a)-P(a-1,a-1)=\Delta_{a-1,a-1}+\Delta_{a,a-1}
 \ge H^P_{a-1}+\Delta_{a,a}\ge H^P_{a-1}+H^P_a .
\]
Substituting $P(a,a)=\frac{3a-1}2H^P_a$ and rearranging yields
$H^P_a\ge\bigl(1+\frac1{3(a-1)}\bigr)H^P_{a-1}$, hence $H^P_a\ge H_a$
for all $a$; this is the source's normalized-growth bound. Finally,
for $b\ge a$, \eqref{eq:slope} gives
\[
 P(a,b)\ \ge\ P(a,a)\Bigl(1+\frac{2(b-a)}{3a-1}\Bigr)
 =H^P_a\Bigl(b+\frac{a-1}2\Bigr)\ \ge\ H_a\Bigl(b+\frac{a-1}2\Bigr)
 =P_{\min}(a,b),
\]
and symmetry covers $b<a$.

\emph{Comparison with $P_*$.} For $m\ge1$,
$(3m+2)\cdot27m^3-(3m+1)^3(3m-1)=6m+1>0$, so
$(1+\frac1{3m})^3<\frac{3m+2}{3m-1}$. Multiplying over $m<u$ and
telescoping, $H_u^3<\frac{3u-1}2$ for $u\ge2$, and the quotient
$\frac{3u-1}{2}/H_u^3$ is a product of factors exceeding $1$, hence
strictly increasing in $u$. Its limit follows from
$\Gamma(u+\tfrac13)/\Gamma(u)\sim u^{1/3}$, which also gives the
diagonal asymptotics. The envelope bound is inherited from $P_*$.
\end{proof}

\begin{corollary}\label{cor:sharp}
Among all profiles satisfying the source's hypotheses (symmetry,
$P(1,b)=b$, discrete concavity, shifted tripling), the minimum possible
diagonal value is exactly
\[
 \min P(a,a)=\frac{3a-1}{2}\prod_{m=1}^{a-1}\Bigl(1+\frac1{3m}\Bigr)
 \qquad(a\ge1),
\]
attained by $P_{\min}$. Thus the source's diagonal-growth lemma is sharp
at every $a$, not only in the exponent $4/3$: its intermediate bound
$H_a\ge\prod_{m<a}(1+1/(3m))$ cannot be improved, and the best constant
$c$ with $P(a,a)\ge c\,a^{4/3}(1+o(1))$ for all admissible $P$ is
$c=3/(2\Gamma(4/3))$.
\end{corollary}

Consequently the barrier is exact in the following sense. Any proof
that uses only constraints satisfied by $P_{\min}$ cannot show $t<3/4$,
since $P_{\min}\le(a+b-1)^{4/3}$. Because $P_{\min}$ is the least
admissible profile, it is the canonical first test for any proposed new
inequality: if $P_{\min}$ satisfies the inequality, adding it to the
system leaves the barrier intact. The earlier profile $P_*$ is not
least; it exceeds $P_{\min}$ by a factor below $1.0223$ everywhere.

\paragraph{Sectors for the least profile.}
Lemma~\ref{lem:ratio} also reads
$P_{\min}(a,h)\le H_a(h+d)$ with $d=(a-1)/2$, with equality if and only
if $h\ge a-1$. Repeating the argument of
Section~\ref{sec:sector-barrier} verbatim with $L$ replaced by $H_a$,
the least profile satisfies every symmetrized matched-sector
consequence $P(a,B)\ge\sum_iP(a,h_i)$ under the packing
bound~\eqref{eq:packing-bound}, and for $s\ge2$, $B\ge B_{\min}$,
\[
 P_{\min}(a,B)-\sum_iP_{\min}(a,h_i)
 =H_a(B-B_{\min})+H_a(2\lfloor s/2\rfloor-s+1)d
  +\sum_i\bigl[H_a(h_i+d)-P_{\min}(a,h_i)\bigr].
\]
For $a>1$ the equality cases differ from those of $P_*$ by one unit:
odd $s$ saturates exactly when $B=B_{\min}$ and every $h_i\ge a-1$
(rather than $h_i\ge a$), and even $s$ has gap at least $H_ad$, which
is smaller than the slack $Ld$ of $P_*$. For $a=1$ nothing changes.
As before, these are identities of an abstract profile and say nothing
about realizability by tensor degenerations.

\section{Before symmetrization: the leg-resolved system}\label{sec:leg}
The source derives its inequalities for each of the six leg-permuted
characters and only then multiplies them into the symmetrized profile.
This section records what the inequalities say \emph{before} that
product is taken, and shows that the barrier survives there too.

Number the legs of $C(a,b)$ as $1$ (first input, dimension $a$),
$2$ (second input, dimension $b$) and $3$ (output dual, dimension
$a+b-1$). A \emph{placement} is a bijection $\pi$ from $\{1,2,3\}$ to
the legs $\{X,Y,Z\}$ of a character $\lambda$ with dot-product
exponents $(p_X,p_Y,p_Z)\in[0,1]^3$; write $C^\pi$ for $C$ with leg
$j$ placed on $\pi(j)$ and
\[
 \ell_\pi(a,b)=\lambda\bigl(C(a,b)^\pi\bigr)>0 .
\]
Let $\tau$ exchange $1,2$ and $\sigma$ exchange $2,3$. Since $C(b,a)$ is
$C(a,b)$ with its input legs exchanged, $\ell_\pi(b,a)=\ell_{\pi\tau}(a,b)$.
For $p>0$ and $A_1,\dots,A_k>0$ put
$\Phi_p(A_1,\dots,A_k)=\bigl(\sum_iA_i^{1/p}\bigr)^p$, and
$\Phi_0(A_1,\dots,A_k)=\max_iA_i$.

\begin{lemma}[H\"older form]\label{lem:hoelder}
For $p>0$,
$\sup_q e^{pH(q)}\prod_iA_i^{q_i}=\Phi_p(A_1,\dots,A_k)$, the supremum
over probability vectors being attained at $q_i=A_i^{1/p}/\sum_jA_j^{1/p}$.
For $p=0$ the supremum is $\max_iA_i$.
\end{lemma}
\begin{proof}
The objective $\sum_iq_i\log A_i-p\sum_iq_i\log q_i$ is concave in $q$.
At the stated $q$, the partial derivatives
$\log A_i-p\log q_i-p$ all equal $p\log\sum_jA_j^{1/p}-p$, so this
interior point is the maximizer, and the value there is
$\sum_iq_i\bigl(\log A_i-p\log q_i\bigr)=p\log\sum_jA_j^{1/p}$.
The case $p=0$ is immediate.
\end{proof}

Applying Lemma~\ref{lem:hoelder} to the source's shared-leg entropy
inequality for the character $\lambda_\pi$, with $p=p_{\pi(1)}$ the
exponent of the leg carrying the shared first input, the determinant
filtration and the three-sector degeneration give, for every placement,
\begin{align}
 \ell_\pi(1,b)&=b^{\,p_{\pi(1)}},\qquad
 \ell_\pi(a,1)=a^{\,p_{\pi(2)}},\qquad
 1\le\ell_\pi(a,b)\le a+b-1,\label{eq:L1}\\
 2^{p}\,\ell_\pi(a,b)&\ \ge\ \Phi_p\bigl(\ell_\pi(a,b+1),\ell_\pi(a,b-1)\bigr)
 \qquad(b\ge2),\label{eq:L3}\\
 \ell_\pi(a,3h+a-1)&\ \ge\ \Phi_p\bigl(\ell_\pi(a,h),\ell_\pi(a,h),
 \ell_{\pi\sigma}(a,h)\bigr),\label{eq:L4}
\end{align}
together with monotonicity in each argument and, for matched sectors
of lengths $h_i$ with orientations $\varepsilon_i\in\{\mathrm{id},\sigma\}$
as in Section~\ref{sec:sector-barrier},
$\ell_\pi(a,B)\ge\Phi_p\bigl(\ell_{\pi\varepsilon_1}(a,h_1),\dots,
\ell_{\pi\varepsilon_s}(a,h_s)\bigr)$. The boundary values come from
$C(1,b)=x_0\sum_jy_jz_j$ and $C(a,1)=y_0\sum_ix_iz_i$, the dot products
whose scalar leg is leg $1$, respectively leg $2$; the first-argument
versions of
\eqref{eq:L3}--\eqref{eq:L4} are the same statements for $\pi\tau$.
We call a family $(\ell_\pi)_\pi$ with exponents $p$ satisfying all
of these a \emph{leg-resolved solution} with $3t=p_X+p_Y+p_Z$.
Multiplying \eqref{eq:L3} over the six placements and using
$\sum_\pi p_{\pi(1)}=6t$ recovers the source's concavity for
$P=(\prod_\pi\ell_\pi)^{1/(6t)}$, and likewise for tripling and
sectors; so every symmetrized inequality is implied, and a barrier for
the leg-resolved system is stronger than one for the profile.

\begin{theorem}\label{thm:leg}
In the coordinates $(\log\ell_\pi,p)$ the leg-resolved solutions form a
convex set. It contains
\begin{enumerate}
\item the three flattening characters: flattening along the leg $F$
has $p_F=0$ and the other two exponents $1$, and
$\ell_\pi(a,b)\in\{a,\,b,\,a+b-1\}$ according to whether leg $1$, $2$
or $3$ of $C$ is placed on $F$; here $3t=2$;
\item the symmetric assignment $\ell_\pi\equiv P_{\min}^{3/4}$ for all
$\pi$, $p=(\tfrac34,\tfrac34,\tfrac34)$; here $3t=9/4$.
\end{enumerate}
Hence for every $3t\in[2,9/4]$ and every exponent triple in the convex
hull of $(0,1,1),(1,0,1),(1,1,0),(\tfrac34,\tfrac34,\tfrac34)$ there is
a leg-resolved solution, and none of the inequalities above, used with
arbitrary probability vectors and arbitrary matched sectors but without
symmetrization, excludes $3t=9/4$.
\end{theorem}
\begin{proof}
For each fixed probability vector $q$, each inequality in its entropy
form reads
$\log\ell_\pi(a,b)+p\log2\ge pH(q)+\sum_iq_i\log\ell_{\pi_i}(\cdot)$,
which is linear in $(\log\ell,p)$; the H\"older forms are the
intersections over $q$. The boundary conditions, the rank and
monotonicity constraints and the exponent constraints are linear as
well, so the solution set is convex.

Flattening ranks are characters, so they satisfy every inequality of
the source. The listed values are their flattening ranks on
$C(a,b)$: the leg-$1$ and leg-$2$ flattenings have full rank $a$ and $b$
because their rows have disjoint supports, and the output flattening
has rank $a+b-1$. The exponents follow from the flattening ranks of the
oriented dot products. The tests verify \eqref{eq:L1}--\eqref{eq:L4}
directly for all six placements, where \eqref{eq:L3} and \eqref{eq:L4}
hold with equality.

For the symmetric assignment, \eqref{eq:L1} is $P_{\min}(1,b)=b$,
$P_{\min}(a,1)=a$ and $P_{\min}\le(a+b-1)^{4/3}$; raising
\eqref{eq:L3} and \eqref{eq:L4} to the power $4/3$ turns them into the
discrete concavity and shifted tripling of $P_{\min}$, and the sector
inequalities into those of the preceding section. Symmetry of $P_{\min}$
supplies the placements $\pi\tau$. Theorem~\ref{thm:least} and the
sector identities give all of these. The remaining assertions follow
from convexity.
\end{proof}

\begin{corollary}[An exponent constraint]\label{cor:exponent}
Let $\lambda$ be any tensor character. If two of $p_X,p_Y,p_Z$ equal
$1$, the third equals $0$, and then $\lambda(T_m)=m^2$ for all $m$.
\end{corollary}
\begin{proof}
Say $p_X=p_Y=1$ and take the identity placement, so $p=1$ and
$\Phi_1$ is the sum. Inequality~\eqref{eq:L4} with $h=1$ and the rank
bound in \eqref{eq:L1} give
\[
 2a+1\ \ge\ \ell_{\mathrm{id}}(a,a+2)\ \ge\
 2\ell_{\mathrm{id}}(a,1)+\ell_{\sigma}(a,1)=2a+a^{p_Z},
\]
using $\ell_{\mathrm{id}}(a,1)=a^{p_Y}$ and $\ell_\sigma(a,1)=a^{p_Z}$.
Thus $a^{p_Z}\le1$ for every $a$, so $p_Z=0$, and
$\lambda(T_m)=m^{p_X+p_Y+p_Z}=m^2$.
\end{proof}

The corollary shows that the leg-resolved system does retain
information that the symmetrized profile discards: on the slice
$p_X+p_Y+p_Z=9/4$ the triples $(1,1,\tfrac14)$ and its permutations are
impossible, and the same inequality at general $(a,h)$ with the
monotone lower bounds $\ell_\pi(a,h)\ge\max(a^{p_{\pi(2)}},h^{p})$
excludes a neighbourhood of each of these three corners
(\texttt{progress/exponent\_region.py} maps it). That information
constrains the \emph{shape} of the exponent triple, not its sum: by
Theorem~\ref{thm:leg} the symmetric triple with sum $9/4$ remains
untouched. Whether asymmetric triples with sum $9/4$ admit leg-resolved
solutions is left open; it does not affect the barrier.

\paragraph{Reproducibility.}
The accompanying README documents \texttt{python -B run\_checks.py}.
The scripts, source manifest, and test expectations are in this repository.
The LaTeX source was compiled successfully with the Codex document editor.
\end{document}
